\documentclass[french,12pt,a4paper]{book}
\usepackage{gensymb} 
\usepackage[utf8]{inputenc}
\usepackage[T1]{fontenc}
\usepackage[english]{babel}
\usepackage[default, scale=.95]{opensans} % police Open Sans
\usepackage{fontawesome}
\usepackage{amsmath}
\usepackage{amsfonts}
\usepackage{fancyhdr}
\usepackage{amssymb}
\usepackage{xcolor} % où color selon l'installation
\definecolor{Prune}{RGB}{99,0,60} % l14-33 : couleurs de la charte graphique upsaclay
\definecolor{B1}{RGB}{49,62,72} 
\definecolor{C1}{RGB}{124,135,143}
\definecolor{D1}{RGB}{213,218,223}
\definecolor{A2}{RGB}{198,11,70}
\definecolor{B2}{RGB}{237,20,91}
\definecolor{C2}{RGB}{238,52,35}
\definecolor{D2}{RGB}{243,115,32}
\definecolor{A3}{RGB}{124,42,144}
\definecolor{B3}{RGB}{125,106,175}
\definecolor{C3}{RGB}{198,103,29}
\definecolor{D3}{RGB}{254,188,24}
\definecolor{A4}{RGB}{0,78,125}
\definecolor{B4}{RGB}{14,135,201}
\definecolor{C4}{RGB}{0,148,181}
\definecolor{D4}{RGB}{70,195,210}
\definecolor{A5}{RGB}{0,128,122}
\definecolor{B5}{RGB}{64,183,105}
\definecolor{C5}{RGB}{140,198,62}
\definecolor{D5}{RGB}{213,223,61}
\usepackage{mdframed}
\usepackage{multirow} %% Pour mettre un texte sur plusieurs rangées
\usepackage{multicol} %% Pour mettre un texte sur plusieurs colonnes
\usepackage{scrextend} %Forcer la 4ème  de couverture en page pair
\usepackage{tikz}
\usepackage{graphicx}
\usepackage[absolute]{textpos} 
\usepackage{colortbl}
\usepackage{array}
\usepackage{geometry}
\usepackage{titlesec}
\usepackage{tocloft}
\usepackage{hyperref}
\usepackage[toc]{glossaries}
\makeglossaries

\newacronym{epr}{EPR}{Einstein, Podolsky and Rosen}
\newacronym{chsh}{CHSH}{Clauser-Horne-Shimony-Holt}
\newacronym{sos}{SOS}{Sum-of-Squares}
\newacronym{CPTP}{CPTP}{completely-positive trace-preserving}

\newacronym{povm}{POVM}{positive-value operator measurement}
\newacronym{DIQKD}{DIQKD}{Device-Independent Quantum Key Distribution}
\newacronym{QKD}{QKD}{Quantum Key Distribution}

\newacronym{rl}{RL}{Reinforcement Learning}

\hypersetup{ % paramétrage couleur des liens hypertextes, toujours garder colorlinks=true
	colorlinks=true,
	linkcolor=black,
	urlcolor=purple}
 
\makeatletter
\renewcommand{\numberline}[1]{%
  \@cftbsnum #1\@cftasnum~\@cftasnumb%
}
\makeatother

\pagestyle{plain} % pour ne garder que les n°de page en milieu-bas et supprimer les indications de chapitre en marge haute

\usepackage{lipsum} % à retirer!!!
\begin{document}
	
	\begin{titlepage}
		
		%\thispagestyle{empty}
		
		\newgeometry{left=6cm,bottom=2cm, top=1cm, right=1cm}
		
		\tikz[remember picture,overlay] \node[opacity=1,inner sep=0pt] at (-13mm,-135mm){\includegraphics{Logo/Frame-ups.pdf}};
		
		%*****************************************************
		%******** NUMÉRO D'ORDRE DE LA THÈSE À COMPLÉTER *****
		%******** POUR LE SECOND DÉPOT                   *****
		%*****************************************************
		
		\color{white}
		
		\begin{picture}(0,0)
			\put(-152,-743){\rotatebox{90}{\Large \textsc{THESE DE DOCTORAT}}} \\
			\put(-120,-743){\rotatebox{90}{NNT : 2020UPASA001}}
		\end{picture}
		
		%*****************************************************
		%**  LOGO  ÉTABLISSEMENT PARTENAIRE SI COTUTELLE
		%**  CHANGER L'IMAGE PAR DÉFAUT **
		%*****************************************************
		\vspace{-14mm} % à ajuster en fonction de la hauteur du logo
		%\flushright \includegraphics[scale=1]{logo2.png}
		
		%*****************************************************
		%******************** TITRE **************************
		%*****************************************************
		
		\flushright
		\vspace{30mm} % à régler éventuellement
		\color{Prune}
		\fontfamily{cmss}\fontseries{m}\fontsize{22}{26}\selectfont
		\Huge Device-independent certification: from foundations to applications \\
		
		\normalsize
		\color{black}
		\Large{\textit{Device-independent certification: Fondations et Applications}}\\
		%*****************************************************
		
		\fontfamily{fvs}\fontseries{m}\fontsize{8}{12}\selectfont
		
		\vspace{1.5cm}
		
		\normalsize
		\textbf{Thèse de doctorat de l'université Paris-Saclay} \\
		
		\vspace{6mm}
		
		\small École doctorale n\degree564 : physique en Île-de-France (PIF)\\
		\small Spécialité de doctorat: Physique\\
		\small Graduate School : Physique, Référent : Faculté des sciences d'Orsay \\
		\vspace{6mm}
		
		\footnotesize Thèse préparée à l’Institut de Physique Théorique (Université Paris-Saclay, CNRS, CEA) sous la direction de Nicolas Sangouard, directeur de recherche (CEA). \\
		\vspace{15mm}
		
		\textbf{Thèse soutenue à Paris-Saclay, le JJ mois AAAA, par}\\
		\bigskip
		\Large {\color{Prune} \textbf{Xavier VALCARCE}}
		
		%************************************
		\vspace{\fill} % ALIGNER LE TABLEAU EN BAS DE PAGE
		%************************************
		
		\bigskip
		
		\flushleft
		\small \textbf{Composition du jury}\\
		\vspace{2mm}
		\scriptsize
		\begin{tabular}{|p{7cm}l}
			\arrayrulecolor{Prune}
			\textbf{Prénom Nom} &   Président ou Présidente\\ 
			Titre, Affiliation & \\
			\textbf{Prénom Nom} &  Rapporteur \& Examinateur / trice \\ 
			Titre, Affiliation   &   \\ 
			\textbf{Prénom Nom} &  Rapporteur \& Examinateur / trice \\ 
			Titre, Affiliation  &   \\ 
			\textbf{Prénom Nom} &  Examinateur ou Examinatrice \\ 
			Titre, Affiliation   &   \\ 
			\textbf{Prénom Nom} &  Examinateur ou Examinatrice \\ 
			Titre, Affiliation   &   \\ 
			\textbf{Nicolas Sangouard} &  Directeur ou Directrice de thèse \\ 
			Titre, Affiliation   &   \\ 
			
		\end{tabular} 
		
	\end{titlepage}

\Ifthispageodd{\newpage\thispagestyle{empty}\null\newpage}{}
\thispagestyle{empty}
\newgeometry{top=1.5cm, bottom=1.25cm, left=2cm, right=2cm}
\fontfamily{rm}\selectfont

\lhead{}
\rhead{}
\rfoot{}
\cfoot{}
\lfoot{}

\noindent 
%*****************************************************
%***** LOGO DE L'ED À CHANGER IMPÉRATIVEMENT *********
%*****************************************************
\includegraphics[height=2.45cm]{Logo/logo_usp_PIF.png}
\vspace{1cm}
%*****************************************************
\fontfamily{cmss}\fontseries{m}\selectfont

\small

\begin{mdframed}[linecolor=Prune,linewidth=1]
	
	\textbf{Titre:} titre (en français)..................................................................................................................
	
	\noindent \textbf{Mots clés:} 3 à 6 mots clefs (version en français)
	
	\vspace{-.5cm}
	\begin{multicols}{2}
		\noindent \textbf{Résumé:}\lipsum[1-2] 
	\end{multicols}
	
\end{mdframed}

\vspace{8mm}

\begin{mdframed}[linecolor=Prune,linewidth=1]
	
	\textbf{Title:} titre (en anglais)..................................................................................................................
	
	\noindent \textbf{Keywords:} 3 à 6 mots clefs (version en anglais)
	
	\begin{multicols}{2}
		\noindent \textbf{Abstract:} \lipsum[1-2]
	\end{multicols}
\end{mdframed}

\titleformat{\part}[block]{\bfseries\Huge\color{Prune}\filright}{\thepart\ --\ }{.1ex}
{\vspace{0.1ex}
}
[\vspace{1ex}]

\titleformat{\chapter}[hang]{\bfseries\Large\color{Prune}}{\thechapter\ -}{.1ex}
{\vspace{0.1ex}
}
[\vspace{1ex}]
\titlespacing{\chapter}{0pc}{0ex}{0.5pc}

\titleformat{\section}[hang]{\bfseries\normalsize}{\thesection\ .}{0.5pt}
{\vspace{0.1ex}
}
[\vspace{0.1ex}]
\titlespacing{\section}{1.5pc}{4ex plus .1ex minus .2ex}{.8pc}

\titleformat{\subsection}[hang]{\bfseries\small}{\thesubsection\ .}{1pt}
{\vspace{0.1ex}
}
[\vspace{0.1ex}]
\titlespacing{\subsection}{3pc}{2ex plus .1ex minus .2ex}{.1pc}

\newcounter{article}
\newcommand{\article}[1]{%
  \par\refstepcounter{article}% Increase section counter
  \sectionmark{#1}% Add section mark (header)
  \addcontentsline{toc}{chapter}{\protect\numberline{Article \thearticle}#1}% Add section to ToC
  % Add more content here, if needed.
}
\frontmatter

\newgeometry{top=4cm, bottom=4cm, left=2cm, right=2cm}

\chapter{Acknowledgements}
\newpage

\tableofcontents

\newgeometry{top=4cm, bottom=4cm, left=4cm, right=4cm}

\printglossary[type=\acronymtype,title=Acronyms, toctitle=Acronyms]

\chapter{Thesis overview}

\section{Motivations and background}

\section{Summary of the contributions}

\chapter{Synthèse en français}

\section{Motivations et context}

\section{Résumé des contributions}

\chapter{Organisation of the manuscript}

\chapter{Publications}

\mainmatter

\part{Self-testing for device-independent certification}

\chapter{Overview and definitions}

\chapter{Self-testing of bipartite states}

\chapter{Robust self-testing}

\chapter{Self-testing as a device-independent certificate}

\section{Device-independent randomness generation}

\section{Device-independent quantum cryptography}

\titleformat{\article}
   {\normalfont\Large\bfseries}{\thearticle}{1em}{}

\chapter*{Article 1}
\label{article_minCHSH}
\article{What is the minimum CHSH score certifying that a state resembles the singlet?}

\begin{center}
\textrm{\LARGE What is the minimum CHSH score certifying that a state resembles the singlet?}

\vspace{2cm}

\normalsize
Xavier Valcarce$^1$, Pavel Sekatski$^1$, Davide Orsucci$^1$, Enky Oudot$^{1,2}$, Jean-Daniel Bancal$^{1,2}$, and Nicolas Sangouard$^1$
\bigbreak

{\footnotesize
$^1$ Departement Physik, Universität Basel, Klingelbergstraße 82, 4056 Basel, Schweiz \\
$^2$ Département de Physique Appliquée, Université de Genève, 1211 Genève, Suisse
}

\raggedright
\bigbreak
\faLink \quad \href{https://quantum-journal.org/papers/q-2020-03-23-246/}{Quantum, volume 4, page 246} \\
\faLink \quad \href{https://arxiv.org/abs/1910.04606}{arXiv preprint: 1910.04606}
\vspace{1cm}

\centering
\textbf{Abstract}
\bigbreak

A quantum state can be characterized from the violation of a Bell inequality.
The well-known CHSH inequality for example can be used to quantify the fidelity (up to local isometries) of the measured state with respect to the singlet state.
In this work, we look for the minimum CHSH violation leading to a non-trivial fidelity.
In particular, we provide a new analytical approach to explore this problem in a device-independent framework, where the fidelity bound holds without assumption about the internal working of devices used in the CHSH test.
We give an example which pushes the minimum CHSH threshold from $\approx 2.0014$ to $\approx 2.05$, far from the local bound.
This is in sharp contrast with the device-dependent (two-qubit) case, where entanglement is one-to-one related to a non-trivial singlet fidelity.
We discuss this result in a broad context including device-dependent/independent state characterizations with various classical resources.

\end{center}

\chapter*{Article 2}
\article{Self-testing two-qubit maximally entangled states from generalized Clauser-Horne-Shimony-Holt tests}

\centering
\textrm{\LARGE Self-testing two-qubit maximally entangled states from generalized Clauser-Horne-Shimony-Holt tests}

\vspace{2cm}

\normalsize
Xavier Valcarce$^1,2$, Julian Zivy$^1,2$, Nicolas Sangouard$^1,2$, and Pavel Sekatski$^2$
\bigbreak

{\footnotesize
	$^1$ Université Paris-Saclay, CEA, CNRS, Institut de Physique Théorique, 91191 Gif-sur-Yvette, France \\
	$^2$ Departement Physik, Universität Basel, Klingelbergstraße 82, 4056 Basel, Schweiz
}

\raggedright
\bigbreak
\faLink \quad \href{https://journals.aps.org/prresearch/abstract/10.1103/PhysRevResearch.4.013049}{Phys. Rev. Research 4, 013049} \\
\faLink \quad \href{https://arxiv.org/abs/2011.03047}{arXiv preprint: 2011.03047}
\vspace{1cm}

\centering
\textbf{Abstract}
\bigbreak

Device-independent certification, also known as self-testing, aims at guaranteeing the proper functioning of untrusted and uncharacterized devices.
For example, the quality of an unknown source expected to produce two-qubit maximally entangled states can be evaluated in a bi-partite scenario, each party using two binary measurements.
The most robust approach consists in deducing the fidelity of produced states with respect to a two-qubit maximally entangled state from the violation of the CHSH inequality.
In this paper, we show how the self-testing of two-qubit maximally entangled states is improved by a refined analysis of measurement statistics.
The use of suitably chosen Bell tests, depending on the observed correlations, allows one to conclude higher fidelities than ones previously known.
In particular, nontrivial self-testing statements can be obtained from correlations that cannot be exploited by a CHSH-based self-testing strategy.
Our results not only provide novel insight into the set of quantum correlations suited for self-testing, but also facilitate the experimental implementations of device-independent certifications.


\part{Device-independent communication and cryptography}

\chapter{Device-independent quantum key distribution}

\section{Protocols overview}

\section{Entropy bound}

\section{Finite-size analysis}

\chapter{Implementing device independent quantum key distribution}

\section{Plateform of choice: Photonic devices}

\section{Automated design of photonic experiment}

\chapter*{Article 3}
\article{Device-independent quantum key distribution from generalized CHSH inequalities}

\centering
\textrm{\LARGE Device-independent quantum key distribution from generalized CHSH inequalities}

\vspace{2cm}

\normalsize
Pavel Sekatski$^{1}$, Jean-Daniel Bancal$^{2}$, Xavier Valcarce$^{3}$, Ernest Y.-Z. Tan$^{4}$, Renato Renner$^{4}$, and Nicolas Sangouard$^{1,3}$
\bigbreak

{\footnotesize
	$^1$ Departement Physik, Universität Basel, Klingelbergstraße 82, 4056 Basel, Schweiz \\
	$^2$ Department of Applied Physics, University of Geneva, Chemin de Pinchat 22, 1211 Geneva, Switzerland \\
	$^3$ Université Paris-Saclay, CEA, CNRS, Institut de Physique Théorique, 91191 Gif-sur-Yvette, France \\
	$^4$ Institute for Theoretical Physics, ETH Zürich, 8093 Zürich, Switzerland
}

\raggedright
\bigbreak
\faLink \quad \href{https://quantum-journal.org/papers/q-2021-04-26-444/}{Quantum, volume 5, page 444} \\
\faLink \quad \href{https://arxiv.org/abs/2009.01784v4}{arXiv preprint: 2009.01784v4}
\vspace{1cm}

\centering
\textbf{Abstract}
\bigbreak

Device-independent quantum key distribution aims at providing security guarantees even when using largely uncharacterised devices.
In the simplest scenario, these guarantees are derived from the CHSH score, which is a simple linear combination of four correlation functions.
We here derive a security proof from a generalisation of the CHSH score, which effectively takes into account the individual values of two correlation functions.
We show that this additional information, which is anyway available in practice, allows one to get higher key rates than with the CHSH score.
We discuss the potential advantage of this technique for realistic photonic implementations of device-independent quantum key distribution.

\end{center}

\chapter*{Article 4}
\article{Automated design of quantum optical experiments for device-independent quantum key distribution}

\begin{center}
\textrm{\LARGE Automated design of quantum optical experiments for device-independent quantum key distribution}

\vspace{2cm}

\normalsize
Xavier Valcarce$^{1}$, Pavel Sekatski$^{2}$, Elie Gouzien$^{1}$, Alexey Melnikov$^{3}$, Nicolas Sangouard$^{1}$
\bigbreak

{\footnotesize
	$^1$ Université Paris-Saclay, CEA, CNRS, Institut de Physique Théorique, 91191 Gif-sur-Yvette, France \\
	$^2$ Department of Applied Physics, University of Geneva, Chemin de Pinchat 22, 1211 Geneva, Switzerland \\
	$^3$ Terra Quantum AG, 9000 St Gallen, Switzerland
}

\raggedright
\bigbreak
\faLink \quad \href{https://arxiv.org/abs/2209.06468}{arXiv preprint: 2209.06468}
\vspace{1cm}

\centering
\textbf{Abstract}
\bigbreak

Device-independent quantum key distribution (DIQKD) reduces the vulnerability to side-channel attacks of standard QKD protocols by removing the need for characterized quantum devices.
The higher security guarantees come however, at the price of a challenging implementation.
Here, we tackle the question of the conception of an experiment for implementing DIQKD with photonic devices.
We introduce a technique combining reinforcement learning, optimisation algorithm and a custom efficient simulation of quantum optics experiments to automate the design of photonic setups maximizing a given function of the measurement statistics.
Applying the algorithm to DIQKD, we get unexpected experimental configurations leading to high key rates and to a high resistance to loss and noise.
These configurations might be helpful to facilitate a first implementation of DIQKD with photonic devices and for future developments targeting improved performances.

\end{center}


\part{Conclusion and perspectives}

\end{document}
